\section{Early Results and Future Plan}

We conducted preliminary experiments on a subset of benchmark programs to evaluate the effectiveness of our branch-aware loop invariant generation approach in addressing the following research question (RQ): ``Can \thistool{} perform better than existing LLM-based loop invariant generation methods in terms of accuracy?'' We also summarize our future plan to further enhance and validate our approach.

\subsection{Preliminary Results}

We selected benchmark programs from the widely used SV-COMP suite, which includes a diverse set of programs with varying complexity and control-flow structures, such as parallel loops and nested loops. First, we implemented state-of-the-art LLM-based loop invariant generation tools, such as Clause2Inv and Lam4Inv, to evaluate whether they could generate correct invariants for these programs. The results show that these tools struggle with programs that have complex control flow, often failing to generate correct invariants within a reasonable time frame.

Next, we extracted the different execution paths from these programs and fed them into the LLMs to generate path-specific invariants. The correctness of each generated invariant was then verified using an SMT solver. We replaced the corresponding paths with these verified invariants and repeated the process until all paths were processed. Finally, we combined all the generated invariants along with the source code as input to the LLM to produce the final invariants for the entire program, which were also checked using the SMT solver.

The preliminary results are promising. Our structure-aware approach successfully generated correct invariants for programs that were previously challenging for existing methods. In particular, we observed a significant improvement in accuracy, with our method achieving a higher success rate in generating correct invariants. This improvement is attributed to the focused reasoning provided by path-specific analysis, which allows the LLM to better understand the context and constraints of each execution path, concentrating on a small portion of the program at a time and thereby reducing the complexity of the reasoning task.

\subsection{Future Directions}

In summary, our future work focuses on expanding the dataset, automating the end-to-end workflow, and improving the effectiveness and scalability of our branch-aware loop invariant generation approach.

First, we plan to extract a more diverse set of benchmark programs from the SV-COMP suite, focusing on programs with complex control structures, intricate arithmetic operations, and array manipulations. This will allow us to better evaluate the generalization capabilities of our approach across different program types and identify areas for improvement.

Second, we aim to implement the algorithms described in Section~\nameref{sec:method} to achieve a fully automated, end-to-end workflow. By integrating path extraction, hierarchical path summarization, LLM-based loop invariant generation, and SMT-based verification, we can minimize manual intervention and streamline the entire loop invariant generation process.

Finally, we plan to conduct a comprehensive evaluation of our approach against state-of-the-art LLM-based loop invariant generation tools on large-scale and industrial benchmarks. This evaluation will consider accuracy, efficiency, and scalability. Additionally, we will explore enhancements such as:
\begin{itemize}
    \item Context-aware LLM prompting strategies to improve invariant precision.
    \item Hybrid techniques combining symbolic reasoning with LLM predictions.
    \item Extending the approach to handle concurrency, recursion, and other advanced constructs.
\end{itemize}

Through these efforts, we aim to develop a robust, scalable, and practical branch-aware loop invariant generation framework that outperforms existing methods.

